
\subsection[\texorpdfstring{$\mathcal M^{\mathrm{main}}_{k+1}(\cdots;\operatorname{top}(\rho))$ 的仓西结构：斯塔谢夫胞腔的 $A_\infty$ 映射对应物}{M main k+1(...;top(rho)) 的仓西结构：斯塔谢夫胞腔的 A infinity 映射对应物}]{%
模空间
$\mathcal M^{\mathrm{main}}_{k+1}
(M',L',\{J_\rho\}_\rho:\beta;\operatorname{top}(\rho))$
的仓西结构：斯塔谢夫胞腔的 $A_\infty$ 映射对应物}
\label{subsec:7-1-5}

现在已在
$\mathcal M^{\mathrm{main}}_{k+1}(\beta)$
上建立了仓西结构。
本书出现的其他种类模空间，
也可以用同样方式构造仓西结构。
不过，我们特别指出：
这里已在博特--莫尔斯框架下
用柱面模型完成了胶合构造。
为在模空间退化时仍得到胶合所需的各种一致估计，
此构造必须谨慎选择函数空间上一族
加权索伯列夫范数。
这种分析并非新的；
规范场论文献 \textbf{[Mro89, DoKr90, Fuk96II]}
及辛几何文献 \textbf{[McSa94, FuOh97]}
早已使用，且在相关专家中已成为成熟的方法。
但这些分析细节在辛几何界似乎尚未广为熟悉，
又不能从上述文献直接借用，
所以为求论述完整并方便读者，
我们在第~7.1.2、7.1.3 小节作了基本自足的说明。

另一方面，就拓扑与几何而言，
要证明这里涉及的各种模空间具有仓西结构，
还有一些新的工作；
这些要点在文献中尚未充分确立，
因此我们详细说明。
具体地说，
我们将证明与定理~\ref{thm:7-1-44} 类似的结果：
当 $M'$ 和 $L'$ 均为一点时，
模空间
$\mathcal M^{\mathrm{main}}_{k+1}
(M',L',\{J_\rho\}_\rho:\beta;\operatorname{top}(\rho))$
是光滑流形。
回忆这个模空间曾用于第~4.6 节
构造滤过 $A_\infty$ 同态，
见定义~\ref{def:4-6-8}。
在该定义中，
对于辛微分同胚 $\psi$，
有
\[
(M',L')=(\psi(M),\psi(L)).
\]
若无须指明 $\psi$ 的定义域和目标，
有时也用 $M,L$ 代替 $M',L'$。

我们指出，
$\mathcal M^{\mathrm{main}}_{k+1}$
可与斯塔谢夫胞腔 \textbf{[Sta63]} 相识别\liftIndex{I233}；
斯塔谢夫曾用它定义
$A_\infty$ 空间的概念，
见定义~7.1.65。
我们还回忆，
$\mathcal M^{\mathrm{main}}_{k+1}(\beta)$
曾用于定义本书滤过 $A_\infty$ 代数的运算
$\mathfrak m_k$。
类似地，下文将证明
$\mathcal N_{k+1}$
同目标空间 $M$ 为一点时的
$\mathcal M^{\mathrm{main}}_{k+1}
(M,L,\{J_\rho\}_\rho:\beta;\operatorname{top}(\rho))$
同胚；
它与 $A_\infty$ 映射的定义有关，
见定义~7.1.66。
我们将在定理~\ref{thm:7-1-51}
中讨论它同斯塔谢夫胞腔的关系。

设
$(\overline C_1,\overline{\mathfrak m})$、
$(\overline C_2,\overline{\mathfrak m}')$
为 $A_\infty$ 代数。
回忆
$\overline{\mathfrak f}_*$
（$*=1,2,\ldots$）
成为 $A_\infty$ 同态的条件可写成
\begin{equation}
\begin{aligned}
&\sum_{\ell,a}
\overline{\mathfrak m}'_\ell
\bigl(
\overline{\mathfrak f}_*
(x_a^{(\ell:1)}),\ldots,
\overline{\mathfrak f}_*
(x_a^{(\ell:\ell)})
\bigr)\\
&\qquad=
\sum_a\pm\overline{\mathfrak f}_*
\bigl(
x_a^{(3:1)},
\overline{\mathfrak m}(x_a^{(3:2)}),
x_a^{(3:3)}
\bigr).
\end{aligned}
\tag{7.1.50}\label{eq:7-1-50}
\end{equation}
其中
$\Delta^{\ell-1}\mathbf x
=\sum_a
x_a^{(\ell:1)}\otimes\cdots\otimes
x_a^{(\ell:\ell)}$。
在下面的定理~\ref{thm:7-1-51}
中，我们将用胞腔分解描述
$\mathcal N_{k+1}$ 的边界；
每个胞腔都是若干
$\mathcal N_{k'+1}$ 或
$\mathcal M^{\mathrm{main}}_{k'+1}$
（$k'<k$）的乘积。
为此先约定记号：
\[
\underline{k}=\{1,\ldots,k\},
\qquad
\underline{k}^{+}
=\underline{k}\cup\{*\},
\]
其中 $*$ 是一个符号。
对每个
$A\in(\underline{k}^{+})^\ell$
且 $\ell<k$，
取 $\mathcal N_{\ell+1}$ 的一个副本，
记为 $\mathcal N_A$；
还取由同一 $A$ 标记的
$\mathcal M^{\mathrm{main}}_{\ell+1}$
的一个副本，
记为 $\mathcal M_A^{\mathrm{main}}$。
对每个 $m$，
另取
$\mathcal M^{\mathrm{main}}_{m+1}$
的一个副本，
记为
$(\mathcal M^{\mathrm{main}}_{m+1})'$。

\begin{theorem}{7.1.51}\label{thm:7-1-51}
对 $k=1,2,\ldots$，
存在微分同胚于 $(k-1)$ 维圆盘
$D^{k-1}$ 的胞腔复形
$\mathcal N_{k+1}$，
其边界分解为下列胞腔之并：
\begin{align}
&\mathcal M^{\mathrm{main}}_{(i+1,\ldots,i+\ell)}
\times
\mathcal N_{(1,\ldots,i,*,i+\ell+1,\ldots,k)}.
\tag{7.1.52.1}\label{eq:7-1-52-1}\\
&\prod_{i=1}^{m}
\mathcal N_{(\ell_{i-1}+1,\ldots,\ell_i)}
\times
(\mathcal M^{\mathrm{main}}_{m+1})'.
\tag{7.1.52.2}\label{eq:7-1-52-2}
\end{align}
这里
$1=\ell_0$、$\ell_m=k$，
且 $\ell_i<\ell_{i+1}-1$。
\translatornote{原书此处的
$\ell_0=1$ 与严格间隔条件，
均无法涵盖证明在印刷页 431 所称的
$m=k$ 单输入分块；
同一证明对 (7.1.52.1) 的端点例也有
标记数不合，见疑误 A-0429。
这里保留原式，不臆定整套索引的改法。}
当目标流形 $M$ 为一点时，
模空间
$\mathcal M^{\mathrm{main}}_{k+1}
(M,L,\{J_\rho\}_\rho:\beta;\operatorname{top}(\rho))$
同胚于上述 $\mathcal N_{k+1}$。
\end{theorem}

式~\eqref{eq:7-1-50} 的右端和左端
分别对应
\eqref{eq:7-1-52-1} 和
\eqref{eq:7-1-52-2}。
也就是说，
$\mathcal M^{\mathrm{main}}_{(i+1,\ldots,i+\ell)}$
对应 $\overline{\mathfrak m}$，
$(\mathcal M^{\mathrm{main}}_{m+1})'$
对应 $\overline{\mathfrak m}'$，
$\mathcal N_{k+1}$
对应 $\overline{\mathfrak f}$。
定理~\ref{thm:7-1-51}
也可以表述为：
$D^{k-1}$ 有
式~\eqref{eq:7-1-52-1}、
\eqref{eq:7-1-52-2}
所述的光滑胞腔分解。

\noindent\textbf{证明。}
先回忆目标空间 $M$ 为一点时
$\mathcal M^{\mathrm{main}}_{k+1}
(M,L,\{J_\rho\}_\rho:\beta;\operatorname{top}(\rho))$
的定义；此时必有 $\beta=0$。
取
\[
(\Sigma,(z_0,\ldots,z_k))
\in\mathcal M^{\mathrm{main}}_{k+1}.
\]
这里 $\Sigma$ 是亏格零有边黎曼曲面，
$z_0,\ldots,z_k$
是遵守次序的边界标记点，
且
$(\Sigma,(z_0,\ldots,z_k))$
稳定；见第~2.1.2 小节。
设
$\Sigma=\bigcup_i\Sigma_i$
为不可约分支分解。
由于没有内部标记点，
稳定性保证不存在球面分支。
按定义~\ref{def:4-6-6}
在不可约分支集上取偏序 $<$。

\begin{definition}{7.1.53}\label{def:7-1-53}
对这样的
$(\Sigma,<)$，
称给每个分支 $\Sigma_i$
赋一个数
$\rho_i\in[0,1]$，
并满足
\[
\Sigma_i\le\Sigma_j
\quad\Longrightarrow\quad
\rho_i\le\rho_j
\]
的赋值 $\rho$ 为
$(\Sigma,<)$ 的\emph{时间分配}\liftIndex{I234}。
\end{definition}

\par\medskip\noindent\begin{minipage}{\textwidth}
\centering
\setlength{\unitlength}{0.65mm}
\begin{picture}(148,83)
\linethickness{0.5pt}
\put(70,39){\circle{27}}
\put(55,58){\circle{19}}
\put(55,72){\circle{10}}
\put(41,58){\circle{9}}
\put(59,17){\circle{18}}
\put(70,39){$\Sigma_0$}
\put(55,58){$\Sigma_1$}
\put(52,71){$\scriptstyle\Sigma_2$}
\put(38,57){$\scriptstyle\Sigma_3$}
\put(59,17){$\Sigma_4$}
\put(84,39){$\scriptstyle z_0$}
\put(80,51){$\scriptstyle z_1$}
\put(66,59){$\scriptstyle z_2$}
\put(62,76){$\scriptstyle z_3$}
\put(54,79){$\scriptstyle z_4$}
\put(45,75){$\scriptstyle z_5$}
\put(29,62){$\scriptstyle z_6$}
\put(28,51){$\scriptstyle z_7$}
\put(43,37){$\scriptstyle z_8$}
\put(44,14){$\scriptstyle z_9$}
\put(60,5){$\scriptstyle z_{10}$}
\put(75,19){$\scriptstyle z_{11}$}
\put(99,66){$\rho_2\le\rho_1$}
\put(99,55){$\rho_3\le\rho_1\le\rho_0$}
\put(99,39){$\rho_4\le\rho_0$}
\end{picture}
\par{\renewcommand{\thefigure}{7.1.19}\refstepcounter{figure}
\normalfont\bfseries 图~\thefigure\label{fig:7-1-19}}
\end{minipage}\par\medskip

见图~\ref{fig:7-1-19}。
第~4.6 节已使用这一概念，
只是当时未为它命名。

\begin{definition}{7.1.54}\label{def:7-1-54}
若 $k\ge2$，
则 $\mathcal N_{k+1}$
是由
$\mathcal M^{\mathrm{main}}_{k+1}$
的稳定元素
$(\Sigma,(z_0,\ldots,z_k))$
及其时间分配 $\rho$
组成的有序对
$((\Sigma,(z_0,\ldots,z_k)),\rho)$
的集合。
见图~\ref{fig:7-1-19}。
依定义，
$\mathcal N_{1+1}$
只有一个点。
\end{definition}

\par\medskip\noindent\begin{minipage}{\textwidth}
\centering
\setlength{\unitlength}{0.75mm}
\begin{picture}(156,86)
\linethickness{0.55pt}
\put(42,69){\line(1,0){7}}
\qbezier(86,65)(91,55)(100,43)
\qbezier(119,64)(119,52)(119,44)
\put(42,24){\line(1,0){7}}
\qbezier(86,24)(92,29)(100,33)
\put(3,58){\framebox(39,22){\shortstack{$\Sigma_{\mathrm m,1}:\ 0$\\$\scriptstyle z_5,z_6$}}}
\put(49,58){\framebox(37,22){\shortstack{$\Sigma_{\mathrm f,2}:\ \frac13$\\$\scriptstyle z_4,z_7$}}}
\put(91,58){\framebox(61,22){\shortstack{$\Sigma_{\mathrm f,3}:\ \frac12,\frac15$\\$\scriptstyle z_1,z_2,z_3$}}}
\put(3,12){\framebox(39,24){\shortstack{$\Sigma_{\mathrm m,0}:\ 0,0,0$\\$\scriptstyle z_8,\ldots,z_{11}$}}}
\put(49,12){\framebox(37,24){\shortstack{$\Sigma_{\mathrm f,1}:\ \frac34$\\$\scriptstyle z_{12}$}}}
\put(100,22){\framebox(48,22){\shortstack{$\Sigma_{\mathrm m'}:\ 1$\\$\scriptstyle z_0,z_{13}$}}}
\end{picture}
\par{\renewcommand{\thefigure}{7.1.20}\refstepcounter{figure}
\normalfont\bfseries 图~\thefigure\label{fig:7-1-20}}
\end{minipage}\par\medskip

利用
$\mathcal M^{\mathrm{main}}_{k+1}$
的拓扑，
可以自然地在 $\mathcal N_{k+1}$
上定义拓扑。
现在证明它满足定理~\ref{thm:7-1-51}
所要求的性质。
为此先引入记号。
写
\[
\mathbf S
=((\Sigma,(z_0,\ldots,z_k)),\rho)
\in\mathcal N_{k+1}.
\]
把所有满足
$\rho_i\in(0,1)$
的不可约分支 $\Sigma_i$
合起来，记为 $\Sigma_{\mathrm f}$；
其连通分支记为
$\Sigma_{\mathrm f,1},\ldots,
\Sigma_{\mathrm f,m}$。
设
$\Sigma_{\mathrm m,1}^{0},\ldots,
\Sigma_{\mathrm m,\ell}^{0}$
为
$\Sigma\setminus\Sigma_{\mathrm f}$
中所有满足 $\rho=0$
的连通分支；
$\ell$ 可以为 $0$。
若存在 $\rho_i=1$ 的不可约分支，
则这些分支的并连通且包含 $z_0$；
将其记为
$\Sigma_{\mathrm m'}^{0}$，
它也可能为空。
分别把
$\Sigma_{\mathrm m,i}^{0}$、
$\Sigma_{\mathrm m'}^{0}$
的闭包记为
$\Sigma_{\mathrm m,i}$、
$\Sigma_{\mathrm m'}$。
把位于这些
$\Sigma_{\mathrm f,i}$、
$\Sigma_{\mathrm m,i}$、
$\Sigma_{\mathrm m'}$
内部的全部奇点平滑化；
参见图~\ref{fig:7-1-20}。
此后忘掉各分支的复结构，
只保留组合型，
记为 $\overline{\mathbf S}$。

$\overline{\mathbf S}$
按如下方式确定一张带状图 $\Gamma$。
图的内部顶点分别对应下列各个分支：
\[
\Sigma_{\mathrm f,1},\ldots,
\Sigma_{\mathrm f,m},
\Sigma_{\mathrm m,1},\ldots,
\Sigma_{\mathrm m,\ell},
\Sigma_{\mathrm m'}
\]
若两个分支相交，
就在相应顶点之间连一条边。
每个标记点对应一个外顶点，
并把它连到包含该标记点的分支所对应的顶点；
这条边称为\emph{外边}。
带状结构由各分支边界上的
标记点及奇点的循环次序确定。
还要记录哪些顶点是外顶点，
以及每个内部顶点的类型
$\mathrm m$、$\mathrm f$ 或 $\mathrm m'$。
由此可把
$\overline{\mathbf S}$
同相应的带状图 $\Gamma$ 相识别；
见图~\ref{fig:7-1-21}。

\par\medskip\noindent\begin{minipage}{\textwidth}
\centering
\setlength{\unitlength}{0.62mm}
\begin{picture}(150,79)
\linethickness{0.55pt}
\put(82,32){\circle*{1.5}}
\put(55,49){\circle*{1.5}}
\put(43,63){\circle*{1.5}}
\put(82,60){\circle*{1.5}}
\put(57,24){\circle*{1.5}}
\put(34,20){\circle*{1.5}}
\qbezier(82,32)(69,40)(55,49)
\qbezier(55,49)(49,56)(43,63)
\qbezier(82,32)(82,46)(82,60)
\qbezier(82,32)(68,28)(57,24)
\qbezier(57,24)(45,22)(34,20)
\qbezier(82,32)(96,32)(108,32)
\qbezier(82,32)(82,20)(82,7)
\qbezier(55,49)(50,44)(44,37)
\qbezier(55,49)(55,59)(55,69)
\qbezier(43,63)(33,63)(24,63)
\qbezier(43,63)(43,70)(43,77)
\qbezier(82,60)(73,69)(66,77)
\qbezier(82,60)(82,70)(82,77)
\qbezier(82,60)(91,69)(99,77)
\qbezier(57,24)(57,16)(57,7)
\qbezier(34,20)(25,29)(18,36)
\qbezier(34,20)(24,19)(15,18)
\qbezier(34,20)(27,11)(20,4)
\qbezier(34,20)(35,10)(35,3)
\put(108,29){$\scriptstyle z_0$}
\put(80,1){$\scriptstyle z_{13}$}
\put(64,70){$\scriptstyle z_1$}
\put(80,72){$\scriptstyle z_2$}
\put(99,72){$\scriptstyle z_3$}
\put(53,70){$\scriptstyle z_4$}
\put(39,75){$\scriptstyle z_5$}
\put(14,60){$\scriptstyle z_6$}
\put(37,35){$\scriptstyle z_7$}
\put(9,34){$\scriptstyle z_8$}
\put(5,14){$\scriptstyle z_9$}
\put(9,0){$\scriptstyle z_{10}$}
\put(32,0){$\scriptstyle z_{11}$}
\put(55,1){$\scriptstyle z_{12}$}
\put(45,65){$\scriptstyle\mathrm m$}
\put(56,52){$\scriptstyle\mathrm f$}
\put(84,62){$\scriptstyle\mathrm f$}
\put(58,26){$\scriptstyle\mathrm f$}
\put(34,23){$\scriptstyle\mathrm m$}
\put(84,35){$\scriptstyle\mathrm m'$}
\end{picture}
\par{\renewcommand{\thefigure}{7.1.21}\refstepcounter{figure}
\normalfont\bfseries 图~\thefigure\label{fig:7-1-21}}
\end{minipage}\par\medskip

对每种这样的组合型 $\Gamma$，
定义
\[
\mathcal N(\Gamma)^\circ
=\{\mathbf S\mid
\overline{\mathbf S}=\Gamma\},
\]
并以
$\mathcal N(\Gamma)$
记其完备化，
也就是它在
$\mathcal N_{k+1}$
中的闭包。
给定 $\Gamma$，
定义
\[
|\Gamma|
=k-1-
\#\bigl\{
v\mid v\text{ 是 }\Gamma
\text{ 的非 }\mathrm f\text{ 型内部顶点}
\bigr\}.
\]
令 $\Gamma_0$
为只有一个内部顶点的图。
对每个 $k$，
依该顶点的类型，
有三个这样的图
$\Gamma_0^{\mathrm f}$、
$\Gamma_0^{\mathrm m}$、
$\Gamma_0^{\mathrm m'}$。
我们对 $k+1$ 归纳证明定理~\ref{thm:7-1-51}。
以下证明假设
对所有 $k'<k$，
该定理已对
$\mathcal N_{k'+1}$
成立。
